% =====================================================================
%  innhold.tex — ALLE FAKTA OM DEG, skrevet én gang.
%
%  Variantfilene i soknader/<firma>/ velger hvilke blokker som brukes,
%  i hvilken rekkefølge, og kan overstyre en blokk med \renewcommand.
%
%  REGEL: fakta endres BARE her. Spissing skjer i variantfilene.
%  Da kan ingen søknad inneholde noe du ikke faktisk har gjort.
%
%  Blokknavnene under (\ErfEn, \UtdTo, \ProsjEn …) er slotter som malen
%  i maler/soknad/cv.tex refererer til. Du kan gi dem navn som beskriver
%  innholdet i stedet (\ErfLagerarbeid, \UtdFagskole) — da må du bytte
%  navnene i malen også.
% =====================================================================

\newcommand{\Navn}{<Fullt navn>}
\newcommand{\NavnKort}{<Slik du signerer>}
% \renewcommand{\Hilsen}{Mvh}                     % valgfri: avslutningen i brev. Standard: Med vennlig hilsen
\newcommand{\Epost}{<e-post>}
\newcommand{\Telefon}{<+47 ...>}
\newcommand{\Sted}{<Sted, fylke>}
\newcommand{\Poststed}{<Poststed>}                  % brukes i datolinja i brev
\newcommand{\NettCV}{}                              % valgfri: adressen til nett-CV-en. Tom = ingen QR-kode
\renewcommand{\Tilgjengelig}{}                      % f.eks. «Kan begynne umiddelbart»
\renewcommand{\Tagline}{}                           % valgfri linje under navnet

% ---------------------------------------------------------------------
%  PROFIL
% ---------------------------------------------------------------------
\newcommand{\Profil}{%
  <Fem til sju linjer om hvem du er. Skriv det én gang, godt. Variantene
  skriver ofte sin egen, men denne er utgangspunktet.>}

\newcommand{\ProfilKort}{%
  <Samme innhold på fire linjer. Brukes på ensides-CV.>}

% ---------------------------------------------------------------------
%  FERDIGHETER — én kommando per gruppe. \sep skiller stikkordene.
%  Ferdigheter er det du kan. Det du har papir på (utdanning, sertifikater,
%  førerkort), har egne blokker og står ikke her.
% ---------------------------------------------------------------------
\newcommand{\FerdEn}{\kompetanse{<Gruppenavn>}{%
  <Stikkord>\sep <Stikkord>\sep <Stikkord>}}
\newcommand{\FerdTo}{\kompetanse{<Gruppenavn>}{%
  <Stikkord>\sep <Stikkord>\sep <Stikkord>}}
\newcommand{\FerdTre}{\kompetanse{<Gruppenavn>}{%
  <Stikkord>\sep <Stikkord>\sep <Stikkord>}}
\newcommand{\FerdFire}{\kompetanse{<Gruppenavn>}{%
  <Stikkord>\sep <Stikkord>\sep <Stikkord>}}

% ---------------------------------------------------------------------
%  ERFARING  —  \oppforing{år}{sted}{rolle}{beskrivelse}
%  Skriv resultater, ikke arbeidsoppgaver. Tall der du har dem.
% ---------------------------------------------------------------------
\newcommand{\ErfEn}{\oppforing{<år–år>}{<Arbeidsgiver>}{<rolle>}{%
  <Hva du gjorde, og hva som ble resultatet.>}}
\newcommand{\ErfTo}{\oppforing{<år–år>}{<Arbeidsgiver>}{<rolle>}{%
  <...>}}
\newcommand{\ErfTre}{\oppforing{<år–år>}{<Arbeidsgiver>}{<rolle>}{%
  <...>}}
\newcommand{\ErfFire}{\oppforing{<år–år>}{<Arbeidsgiver>}{<rolle>}{%
  <...>}}

% Eldre eller mindre relevante jobber, komprimert til én linje.
\newcommand{\EldreKort}{%
  \kortlinje{<år–år>}{\textbf{Annen erfaring:} <jobb> (<år>)\sep <jobb> (<år>)}}

% ---------------------------------------------------------------------
%  UTDANNING OG VERV — to egne seksjoner på CV-en
%  Ikke fullført? Skriv «påbegynt» etter linja, og la datoene si resten.
% ---------------------------------------------------------------------
\newcommand{\UtdEn}{\oppforing{<år–år>}{<Skole>}{<linje/grad>}{%
  <Fag og innhold som er relevant.>}}
\newcommand{\UtdTo}{\oppforing{<år–år>}{<Skole>}{<linje/grad>}{%
  <...>}}
\newcommand{\VervEn}{\oppforing{<år–år>}{<Verv>}{}{%
  <Hva vervet innebar.>}}

\newcommand{\UtdEldreKort}{%
  \kortlinje{<år–år>}{\textbf{Tidligere utdanning:} <skole og linje> (<år>)}}

% ---------------------------------------------------------------------
%  PROSJEKTER  —  \prosjekt{tittel}{beskrivelse}
%  Beskriv HVORFOR prosjektet finnes før HVA det er laget i.
% ---------------------------------------------------------------------
\newcommand{\ProsjEn}{\prosjekt{<Tittel>}{%
  <Hvilket problem løser det, og hva består det av.>}}
\newcommand{\ProsjTo}{\prosjekt{<Tittel>}{%
  <...>}}

% ---------------------------------------------------------------------
%  SPRÅK, SERTIFIKATER, INTERESSER
% ---------------------------------------------------------------------
\newcommand{\Bunn}{\bunn
  {<Språk>\\ <Språk>}
  {<Sertifikat>\\ <Førerkort>}
  {<Interesse>\sep <Interesse>}}

% Kompakt variant for ensides-CV: to linjer i stedet for tre kolonner.
\newcommand{\BunnKort}{\bunnkort
  {<Språk>\sep <Språk>}
  {<Sertifikat>\sep <Førerkort>\sep Referanser oppgis etter forespørsel}}
